\chapter{Supersymmetric Localization On Compact Manifolds}
\label{ch:susy-localization-compact-manifolds}

\ConventionsMixed

Supersymmetric localization is an exact-evaluation method only after the
regulated functional integral has been specified with enough information to
justify deformation of the integration cycle.  This chapter separates the
finite-dimensional theorem from the QFT formulae used in four-dimensional
\(\mathcal N=2\) gauge theory on \(S^4\) and three-dimensional
\(\mathcal N=2\) Chern--Simons--matter theory on \(S^3\).  The formulae below
are therefore not introduced as formal replacements for a definition of a QFT:
they are consequences of a localization datum in the sense of
Definition~\ref{def:susy-localization-datum}, together with the determinant
and instanton-sector prescriptions explicitly stated here.

\section{Finite-Dimensional Deformation Identity}

Let \(X\) be an oriented finite-dimensional supermanifold equipped with an
integration cycle \(C\subset X\) and Berezinian density \(\mu\).  Let \(Q\) be
an odd vector field on \(X\), and let \(B\) be the even vector field \(Q^2\).
For the following proposition all functions and densities are assumed smooth
on a neighborhood of \(C\), and \(C\) is assumed to be either compact without
boundary or equipped with boundary conditions for which the displayed
Stokes term vanishes.

\begin{proposition}[Localization deformation identity]
\label{prop:susy-localization-deformation-identity}
Let \(S,V,\mathcal O\) be functions on \(X\) with parities
\(|S|=|\mathcal O|=0\), \(|V|=1\).  Suppose
\[
  Q S=0,\qquad Q\mathcal O=0,\qquad
  B S=B\mathcal O=BV=0,
\]
and suppose the Berezinian density is \(Q\)-invariant:
\(\mathcal L_Q\mu=0\).  Define
\[
  I(t)=\int_C\mu\,\mathcal O\,\exp\{-S-t\,QV\}.
\]
If \(Q(\mu\,\mathcal O\,V\,\exp\{-S-tQV\})\) has zero integral over \(C\),
then \(I(t)\) is independent of \(t\).
\end{proposition}

\begin{proof}
Because \(Q\) is an odd derivation and \(Q^2V=BV=0\),
\[
  \frac{dI}{dt}
  =
  -\int_C\mu\,\mathcal O\,(QV)\,\exp\{-S-tQV\}.
\]
The hypotheses \(QS=0\), \(Q\mathcal O=0\), and \(Q(QV)=0\) give
\[
\begin{aligned}
  Q\!\left(\mathcal O\,V\,\exp\{-S-tQV\}\right)
  &=
  (Q\mathcal O)V\ee^{-S-tQV}
  +\mathcal O(QV)\ee^{-S-tQV} \\
  &\quad
  -\mathcal O V\,Q(S+tQV)\ee^{-S-tQV} \\
  &=
  \mathcal O(QV)\ee^{-S-tQV}.
\end{aligned}
\]
Using \(\mathcal L_Q\mu=0\), the right hand side integrates to the assumed
zero Stokes term.  Hence \(dI/dt=0\).
\end{proof}

If the real part of \(QV\) is nonnegative on the bosonic body of \(C\), and
if its zero locus \(F\) is compact and clean, the \(t\to\infty\) limit reduces
the integral to \(F\) with a normal Gaussian determinant.  Let \(N_F\) denote
the normal complex to \(F\), obtained by linearizing the \(Q\)-complex in the
directions transverse to \(F\).  In a coordinate chart in which the quadratic
part of \(QV\) is
\[
  (QV)_{\rm quad}
  =
  \frac12 x^i A_{ij}x^j+\psi^a B_{ab}\psi^b+\chi^r C_{ri}x^i ,
\]
the normal contribution is the Berezin Gaussian
\[
  Z_{\rm normal}
  =
  \frac{\operatorname{Pf}(B)\,\det(C)}
  {\sqrt{\det(A)}} ,
\]
with the square-root branch fixed by the original integration cycle.  This
finite formula is the local model for the one-loop determinants in the
compact-space QFT examples below.

\section{QFT Localization Datum}

For a supersymmetric QFT on a compact Riemannian manifold \(M\), the symbol
\[
  Z_M=\int_{\mathfrak F}\mathcal D\Phi\,\ee^{-S[\Phi]}
\]
will be used only as shorthand for a regulated datum:
\[
  (\mathfrak F_\Lambda,C_\Lambda,\mu_\Lambda,S_\Lambda,Q_\Lambda,V_\Lambda)
\]
together with a limiting prescription \(\Lambda\to\infty\).  Here
\(\mathfrak F_\Lambda\) is a finite-mode, lattice, heat-kernel, BV, or other
regulated field-variable space; \(C_\Lambda\) is the declared integration
cycle; \(\mu_\Lambda\) is the regulated Berezinian density; \(Q_\Lambda\) is
an odd symmetry whose square is a combination of an isometry, gauge
transformation, \(R\)-symmetry, and flavor transformation; and \(V_\Lambda\)
is the chosen localizing functional.  The localized expression is justified
when the following conditions are verified or retained as hypotheses:
\[
\begin{gathered}
  Q_\Lambda S_\Lambda=0,\qquad
  \mathcal L_{Q_\Lambda}\mu_\Lambda=0,\qquad
  Q_\Lambda\mathcal O_\Lambda=0,\\
  \int_{\partial C_\Lambda}
  \mu_\Lambda\,\mathcal O_\Lambda V_\Lambda
  \ee^{-S_\Lambda-tQ_\Lambda V_\Lambda}=0,\\
  \operatorname{Re}(Q_\Lambda V_\Lambda)\ge0
  \quad\hbox{on the bosonic integration cycle}.
\end{gathered}
\]
The fixed locus, zero modes, determinants, and singular strata must then be
included in the limiting prescription.
Noncompact Coulomb directions, Jeffrey--Kirwan residues, and point-instanton
compactifications are not automatic outputs of the symbol \(\int\mathcal
D\Phi\).  They are part of the regulated localization datum.  In the formulae
below this datum will be stated by naming the finite-dimensional variables,
the contour, and the determinant or equivariant-integral prescription.

\section{\texorpdfstring{\(S^4\) Localization In Four-Dimensional \(\mathcal N=2\) Gauge Theory}{S4 Localization In Four-Dimensional N=2 Gauge Theory}}

Let \(G\) be compact connected, with Lie algebra \(\mathfrak g\), maximal
torus \(T\), Cartan algebra \(\mathfrak t\), Weyl group \(W_G\), and root
system \(\Delta\subset\mathfrak t^\ast\).  The monograph convention for
non-Abelian gauge theory is
\[
  \operatorname{tr}_{\square}(t^at^b)=\delta^{ab},
  \qquad
  S_{\rm YM}=\frac1{4g_{\rm YM}^2}
  \int \operatorname{tr}_{\square}(F_{\mu\nu}F^{\mu\nu}) .
\]
With \((a,a):=\operatorname{tr}_{\square}(a^2)\) for \(a\in\mathfrak t\), the
round \(S^4_r\) localization variable is the constant adjoint scalar
\(a\in\mathfrak t\) on the Coulomb locus.  The classical Gaussian in this
trace-delta convention is
\[
  Z_{\rm cl}(a)
  =
  \exp\!\left\{
    -\frac{4\pi^2r^2}{g_{\rm YM}^2}(a,a)
  \right\}.
  \label{eq:s4-localization-classical-gaussian}
\]
The coefficient equals the familiar half-trace coefficient because
\(t^a=\sqrt2\,T^a\) and \(g_{\rm ht}^2=2g_{\rm YM}^2\), so
\[
  \frac{8\pi^2}{g_{\rm ht}^2}\operatorname{tr}_{\rm ht}(a_{\rm ht}^2)
  =
  \frac{4\pi^2}{g_{\rm YM}^2}(a,a).
\]

Define the even entire function
\[
  H(x)
  =
  \prod_{n=1}^{\infty}
  \left(1+\frac{x^2}{n^2}\right)^n
  \exp\!\left(-\frac{x^2}{n}\right),
  \label{eq:pestun-H-function}
\]
understood by zeta regularization or by the equivalent Barnes-double-gamma
definition.  For a vector multiplet and hypermultiplets of masses \(m_\ell\)
in representations \(R_\ell\), the perturbative determinant on the round
sphere is
\[
  Z_{1{\rm -loop}}^{S^4}(a)
  =
  \prod_{\alpha\in\Delta_+}
  \alpha(a)^2 H(r\,\alpha(a))^2
  \prod_\ell\prod_{\rho\in{\rm wt}(R_\ell)}
  H\!\left(r\,\rho(a)+r\,m_\ell\right)^{-1}.
  \label{eq:s4-localization-one-loop}
\]
The factor \(\prod_{\alpha>0}\alpha(a)^2\) is the Cartan gauge-fixing
Jacobian.  If the representation is real or pseudoreal, the weights are
counted with the multiplicity appropriate to a full hypermultiplet in the
chosen convention.  A different convention for the hypermultiplet mass shifts
the argument of \(H\) by the corresponding \(R\)-symmetry background; this is
a coordinate change in the localization datum, not a change of the underlying
determinant complex.

\begin{definition}[Nekrasov factor in the \(S^4\) formula]
\label{def:s4-nekrasov-factor}
For the four-dimensional localization formula, \(Z_{\rm inst}(a,\tau;
\epsilon_1,\epsilon_2)\) denotes the equivariant integral over the framed
instanton moduli space on \(\mathbb C^2\), with equivariant rotation weights
\((\epsilon_1,\epsilon_2)\), Coulomb parameter \(a\), and
\[
  \tau=\frac{\theta}{2\pi}+\frac{4\pi\ii}{g_{\rm YM}^2}.
\]
For \(U(N)\) this can be defined by integration over the smooth framed
torsion-free-sheaf compactification.  For other gauge groups and matter
systems the formula requires the corresponding compactification or an
equivalent contour prescription as part of the stated datum.
\end{definition}

The localized partition function is
\[
  Z_{S^4}^{\mathcal N=2}
  =
  \frac1{|W_G|}
  \int_{\mathfrak t} da\,
  Z_{\rm cl}(a)\,
  Z_{1{\rm -loop}}^{S^4}(a)\,
  Z_{\rm inst}\!\left(a,\tau;\frac1r,\frac1r\right)
  Z_{\rm inst}\!\left(a,\bar\tau;\frac1r,\frac1r\right).
  \label{eq:pestun-s4-localization}
\]
The north pole contributes instantons and the south pole contributes
anti-instantons.  The two factors are complex conjugates only when the
integration cycle and masses are chosen so that reflection positivity gives
the corresponding conjugation property.

\begin{example}[\(\mathcal N=4\) Yang--Mills degeneration]
\label{ex:s4-n4-gaussian-matrix-model}
For an \(\mathcal N=2\) vector multiplet plus an adjoint hypermultiplet with
the mass coordinate tuned to the \(\mathcal N=4\) point, the hypermultiplet
weights cancel the \(H\)-factors from the vector multiplet.  The remaining
matrix model is
\[
  Z_{S^4}^{\mathcal N=4}
  =
  \frac1{|W_G|}
  \int_{\mathfrak t}da\,
  \exp\!\left\{
    -\frac{4\pi^2r^2}{g_{\rm YM}^2}(a,a)
  \right\}
  \prod_{\alpha\in\Delta_+}\alpha(a)^2 .
\]
For \(G=U(1)\) there is no root factor and no finite-action instanton sector
on \(\mathbb C^2\).  The partition function reduces to the ordinary Gaussian
\[
  \int_{-\infty}^{\infty} da\,
  \exp\!\left(-\frac{4\pi^2r^2}{g_{\rm YM}^2}a^2\right)
  =
  \frac{g_{\rm YM}}{2\sqrt{\pi}\,r}.
  \label{eq:s4-u1-gaussian-localization}
\]
\end{example}

\section{\texorpdfstring{\(S^3\) Localization In Three-Dimensional \(\mathcal N=2\) Theories}{S3 Localization In Three-Dimensional N=2 Theories}}

Let \(G\), \(T\), \(\mathfrak t\), \(W_G\), and \(\Delta\) be as above.  On
the round \(S^3_r\), the Coulomb-branch localization variable is the constant
real scalar \(\sigma\in\mathfrak t\) in the vector multiplet.  With
Chern--Simons level \(k\), FI parameter \(\zeta\in\mathfrak t^\ast\), chiral
multiplets \(\Phi\) in representations \(R_\Phi\), real masses \(m_\Phi\),
and \(R\)-charges \(\Delta_\Phi\), the Kapustin--Willett--Yaakov matrix
integral is
\[
  Z_{S^3}
  =
  \frac1{|W_G|}
  \int_{\mathfrak t}d\sigma\,
  \ee^{\ii\pi k r^2(\sigma,\sigma)+2\pi\ii r(\zeta,\sigma)}
  \prod_{\alpha\in\Delta_+}
  4\sinh^2\!\bigl(\pi r\,\alpha(\sigma)\bigr)
  \prod_{\Phi}Z_{\Phi}^{S^3}(\sigma),
  \label{eq:kwy-s3-localization}
\]
where
\[
  Z_{\Phi}^{S^3}(\sigma)
  =
  \prod_{\rho\in{\rm wt}(R_\Phi)}
  \exp\!\left[
    \ell\!\left(1-\Delta_\Phi+\ii r\,\rho(\sigma)+\ii r\,m_\Phi\right)
  \right].
\]
The function \(\ell\) is fixed, up to an additive constant that changes the
overall normalization of \(Z_{S^3}\), by
\[
  \ell'(z)=-\pi z\cot(\pi z),
  \qquad
  \ell\!\left(\frac12+\ii x\right)
  +
  \ell\!\left(\frac12-\ii x\right)
  =
  -\log\!\bigl(2\cosh(\pi x)\bigr).
  \label{eq:kwy-ell-identity}
\]
The integration contour is the real Cartan in the displayed unitary
normalization.  If the matter content makes the integral nonconvergent, a
mass, FI parameter, analytic continuation, or contour prescription must be
declared as part of the theory being evaluated.

\begin{example}[Abelian Chern--Simons Gaussian]
\label{ex:s3-u1-cs-fresnel}
For pure \(U(1)_k\) Chern--Simons theory on \(S^3\) with FI parameter
\(\zeta\), the localized integral is the oscillatory Fresnel integral
\[
  Z_{U(1)_k}(\zeta)
  =
  \int_{\mathbb R} d\sigma\,
  \exp\!\left(\ii\pi k\sigma^2+2\pi\ii\zeta\sigma\right).
\]
With the contour obtained by rotating the positive quadratic direction by the
usual \(\ii0\) prescription, completing the square gives
\[
  Z_{U(1)_k}(\zeta)
  =
  \frac{\ee^{\ii\pi\,{\rm sgn}(k)/4}}{\sqrt{|k|}}\,
  \exp\!\left(-\frac{\ii\pi\zeta^2}{k}\right).
  \label{eq:s3-u1-cs-fresnel}
\]
The phase depends on the framing convention of the Chern--Simons theory; the
absolute value and FI dependence follow from the declared contour.
\end{example}

\begin{example}[A conjugate pair of chirals]
\label{ex:s3-chiral-pair-integral}
For two chiral multiplets of charges \(+1\) and \(-1\), both with
\(R\)-charge \(1/2\) and zero real mass, the product of one-loop determinants
on the \(U(1)\) Coulomb parameter is
\[
  \exp\!\left[
    \ell\!\left(\frac12+\ii\sigma\right)
    +
    \ell\!\left(\frac12-\ii\sigma\right)
  \right]
  =
  \frac1{2\cosh(\pi\sigma)}.
\]
The elementary integral
\[
  \int_{-\infty}^{\infty}
  \frac{d\sigma}{2\cosh(\pi\sigma)}
  =
  \frac12
  \label{eq:s3-chiral-pair-integral}
\]
is a finite check on the determinant normalization in
\eqref{eq:kwy-s3-localization}.  It is not by itself a definition of a gauge
theory; the gauge quotient, possible Chern--Simons contact terms, and
background-field couplings remain part of the QFT datum.
\end{example}

\section{Theorem Boundary And Singular Loci}

Equations~\eqref{eq:pestun-s4-localization} and
\eqref{eq:kwy-s3-localization} have different mathematical status from
Proposition~\ref{prop:susy-localization-deformation-identity}.  The
proposition is a finite-dimensional deformation identity under stated
hypotheses.  The compact-space QFT formulae additionally require:
\[
\begin{array}{ll}
S^4:&
\text{a regulator preserving the chosen supercharge, determinant
regularization,}\\
&\text{and an instanton compactification or contour for singular point
instantons;}\\[3pt]
S^3:&
\text{a parity-anomaly regulator, framing convention, real contour or
analytic continuation,}\\
&\text{and a treatment of noncompact Coulomb directions and singular
hyperplanes.}
\end{array}
\]
When these data are supplied, localization turns a protected observable into
an integral over fixed-locus variables.  It does not erase the distinction
between regularization and renormalization, nor does it make contact terms
canonical.  Background Chern--Simons terms in three dimensions and local
curvature counterterms in four dimensions change phases or local
source-dependent factors in the localized answer; they must be tracked as
coordinates of the QFT, not discarded as normalization details.
